\section{Introduction}

RCTX is the shared machine-context contract for the civic evidence package
family. Its job is to say which units exist, in which canonical order, under
which context hash, with which graph identity, source hashes, aligned
attributes, and crosswalks. That contract lets RPLAN, RCOUNT, RHIST, and later
packages refer to the same base dimension without making each package invent a
private geography model.

The boundary is intentionally conservative. RPLAN already has an implemented
\texttt{.rctx} artifact in \texttt{rplan-core} and \texttt{rplan-io}. X.2 does
not replace that path. It treats the current \texttt{rplan\_core::RplanContext}
shape as RCTX-compatible phase-1 infrastructure and keeps full context
construction where it already works. The shared extraction is limited to the
parts that sibling products already need: context unit indexes, source
references, graph records, exact rational crosswalks, package hashing, and
claim-boundary checks.

RMAP is a separate owner concept. It is reserved for rendered or cartographic
presentation: maps, layers, styling, labels, projections for display, tiles,
images, and map-image hashes. RMAP may reference an RCTX context hash, but RCTX
must remain useful without a rendered map. This is the central X.2 claim:
machine context should be shared, while cartographic presentation should not be
mistaken for canonical unit identity.
